% TikZ 图：理想 LC 回路一个周期内的四个典型状态
% Day 24：电磁振荡
\begin{tikzpicture}[scale=0.86,>=Latex,line cap=round,line join=round]
  \path[use as bounding box] (-0.45,-2.20) rectangle (11.15,2.15);

  \foreach \x/\top/\bottom in {
    0/{电场能最大}/{q=+q_m,\ i=0},
    2.85/{磁场能最大}/{q=0,\ i=+I_m},
    5.70/{电场能最大}/{q=-q_m,\ i=0},
    8.55/{磁场能最大}/{q=0,\ i=-I_m}
  }{
    \draw[rounded corners=3pt,fill=gray!5] (\x,0) rectangle ++(2.10,1.30);
    % Closed LC loop: capacitor in the left branch, inductor in the right.
    \draw[thick] (\x+0.50,0.98) -- (\x+1.60,0.98) -- (\x+1.60,0.82);
    \draw[thick,decorate,decoration={coil,aspect=0.45,segment length=4pt,amplitude=3pt}]
      (\x+1.60,0.82) -- (\x+1.60,0.48);
    \draw[thick] (\x+1.60,0.48) -- (\x+1.60,0.32) -- (\x+0.50,0.32);
    \draw[thick] (\x+0.50,0.98) -- (\x+0.50,0.76);
    \draw[thick] (\x+0.36,0.76) -- (\x+0.64,0.76);
    \draw[thick] (\x+0.36,0.54) -- (\x+0.64,0.54);
    \draw[thick] (\x+0.50,0.54) -- (\x+0.50,0.32);
    \node[font=\scriptsize\bfseries] at (\x+1.05,1.62) {\top};
    \node[font=\scriptsize] at (\x+1.05,-0.28) {$\bottom$};
  }

  % Charge and current cues.
  \node[red,font=\scriptsize] at (0.24,0.76) {$+$};
  \node[blue,font=\scriptsize] at (0.24,0.54) {$-$};
  \draw[orange!85!black,->,thick] (3.48,0.98) -- (4.18,0.98)
    node[midway,above,font=\scriptsize] {$i>0$};
  \node[blue,font=\scriptsize] at (5.94,0.76) {$-$};
  \node[red,font=\scriptsize] at (5.94,0.54) {$+$};
  \draw[orange!85!black,->,thick] (10.18,0.98) -- (9.48,0.98)
    node[midway,above,font=\scriptsize] {$i<0$};

  \draw[->,thick] (2.18,0.65) -- (2.72,0.65);
  \draw[->,thick] (5.03,0.65) -- (5.57,0.65);
  \draw[->,thick] (7.88,0.65) -- (8.42,0.65);
  \draw[->,thick] (10.65,-0.72) .. controls (7.40,-1.18) and (3.20,-1.18) ..
    (0.90,-0.72);
  \node[font=\scriptsize,fill=white,inner sep=1.5pt] at (5.55,-1.03)
    {继续演化，回到初态};
  \node[font=\small] at (5.55,-1.78)
    {$T_{\mathrm{LC}}=2\pi\sqrt{LC}$};
\end{tikzpicture}
